\section{Discussion and Conclusion}
\label{sec:conclusion}

This chapter interprets the results, states the contributions against the three
research questions, records the limitations of the study, and lists directions for
future work.

\subsection{Interpretation}
\label{sec:interpretation}

The central result of this thesis is a priority ordering. For a contact-rich
multi-objective planar-pushing task under a single-GPU budget, effort spent on
principled reward design produces a larger and more reliable improvement than
effort spent on structured training dynamics. The three research questions
decompose this result into an efficiency claim, a comparison claim, and a
mechanistic explanation.

The efficiency claim is direct. Potential-based shaping reduces the parallelism
required for a given success rate by a factor of about four, from 2{,}048 to 528
environments. It also improves the combined-objective success rate by 35
percentage points over the earlier reward at the same gate. The reason was
developed in Section~\ref{sec:pbrs_design}. The potential-based reward is dense,
correctly signed on every push, and bounded, whereas the improvement reward is
negative for most useful pushes and has near-zero variance. The policy-invariance
guarantee of Section~\ref{sec:pbrs} allows this dense signal to be added without
changing the task. The result confirms empirically that the guarantee holds under
stochastic contact physics.

The comparison claim is that neither a forced curriculum nor asymmetric self-play
improves on the plain shaped single-agent model. The curriculum is statistically
equivalent to the single-agent model at the large batch used in Isaac Lab, and
self-play is far worse. This does not contradict the literature that motivates
curricula and self-play; it locates a regime in which they do not pay. The
cross-environment result of Section~\ref{sec:cross_env_results} makes the regime
explicit. The curriculum substitutes for the shaped reward when the batch is small
and gradient estimates are noisy, and adds nothing once the batch is large enough
for stable estimates on both objectives. Reward design and curriculum design are
therefore not independent levers of equal standing. At a sufficient batch size,
the shaped reward already supplies what the curriculum would otherwise provide.

The mechanistic explanation concerns why self-play in particular fails.
Section~\ref{sec:rq3} showed that self-play learns the two constituent skills, each
to roughly 50\,\%, but combines them at about 10\,\%, well below the independence
prediction. Three factors act together to produce this outcome, and all are
specific to this task. First, the physics couples translation and rotation, so the
combined gate cannot be satisfied by treating the objectives separately. Second,
the goal distribution is non-stationary, because the goal-proposing agent changes
it at every iteration, which biases the value term of the GAE advantage as shown
in Section~\ref{sec:pbrs_asp}. Third, the combined gate is strict, so partial
competence on each axis does not accumulate into success. The potential-based
reward supplies enough correctly-signed signal to learn each axis, but it cannot
remove the value-term bias that prevents their composition. This is why shaping
lifts self-play above the near-zero level it reaches under the improvement reward,
yet leaves it far below single-agent training.

\subsection{Relation to Prior Work}
\label{sec:relation}

The findings of this thesis agree with the theory of the papers it builds on,
while qualifying the empirical claims of the papers it tests. The distinction
between the two is worth stating precisely.

On reward shaping, the result confirms \cite{Ng1999PolicyInvariance} under
conditions the original theorem did not examine. Ng, Harada, and Russell proved
policy invariance for finite, discrete Markov Decision Processes with
deterministic dynamics. The present task is continuous, high-dimensional, and
driven by stochastic PhysX contact. The shaping potential is a bounded exponential,
rather than the distance-to-goal heuristics of the original paper. The
potential-based reward reaches the same success rate as the hand-tuned reward
while requiring a quarter of the environments, and never exceeds it. This is the
empirical signature the theorem predicts: shaping changes how quickly the optimal
policy is found, not which policy is optimal. The episodic treatment of
\cite{Grzes2017RewardShaping} licences the undiscounted shaping used here, and the
stable value loss in Appendix~\ref{app:curves} matches the bounded-potential
construction. In this respect the thesis extends an established theoretical result
into a contact-rich robotic regime rather than contesting it.

On asymmetric self-play the relationship is more nuanced, and the negative result
must not be overstated. \cite{Sukhbaatar2018IntrinsicMotivation} demonstrated
self-play in reversible or resettable gridworlds and continuous control tasks,
where the goal is a single state to be reached. \cite{OpenAI2021AsymmetricSF}
scaled it to robotic manipulation with large distributed compute and a continuous
action space. The present study differs from both on three axes that
Sections~\ref{sec:pbrs_asp} and \ref{sec:interpretation} identify as decisive. The
goal is multi-objective with a strict combined gate, not a single reachable state.
The action is a discrete macro-action over a large MultiCategorical space, not a
smooth continuous control. The compute budget is a single Graphical Processing
Unit, not a distributed cluster. The finding is therefore not that self-play is
ineffective. Rather, in this specific setting it does not compose the two learned
skills, and it stays far below single-agent training under a budget realistic for
one laboratory. This framing is consistent with \cite{Berner2019Dota2}, whose
success with competitive self-play required months of training at very large
scale. Through that lens, the modest self-play result here supports rather than
contradicts the observation that self-play is demanding of scale.

The behaviour of the Alice Behavioral Cloning term deserves specific comment. It
connects a design choice in \cite{OpenAI2021AsymmetricSF} to an outcome observed
here. Plappert and colleagues introduced Proximal Policy Optimization-style
clipping on the behavioural-cloning loss (Equation~\ref{eq:abc}), to prevent the
imitation term from producing aggressive policy updates. In their continuous-action
grasping domain, the demonstrated actions of the goal-proposing agent stay close to
the goal-solving agent's policy, so the clipping is rarely active and the imitation
signal flows. In the discrete, contact-rich pushing domain of this thesis, the
goal-proposing agent's unguided pushes lie far from the goal-solving agent's
goal-directed policy. The probability ratio in Equation~\ref{eq:abc} is therefore
extreme, the clip saturates, and the imitation gradient vanishes. The same
mechanism that stabilises imitation in the original setting removes it here. This
is not a flaw in the original formulation. It is a consequence of applying that
formulation to a domain whose action distribution differs sharply from the one it
was designed for. This also explains why the two self-play variants differ only
through Alice's reward, not through the cloning term.

On curriculum learning, the regime-dependence of
Section~\ref{sec:cross_env_results} refines rather than rejects the surveys of
\cite{Narvekar2020} and \cite{Portelas2020ACLSurvey}. Those surveys motivate
curricula primarily for settings where reward is sparse or exploration is hard.
The present result is consistent with that motivation and makes it quantitative.
The forced curriculum contributes substantially when the batch is small and
gradient estimates are noisy: at the small \texttt{gym-pusht} batch it raises scene
success from 10.0\,\% to 50.0\,\%. It contributes nothing once the batch is large
enough for stable estimates on both objectives. A well-shaped dense reward and a
curriculum are therefore partial substitutes, not independent and additive
interventions. The point at which one ceases to add value over the other is
governed by the gradient-variance regime. To the best of the author's knowledge,
this observation is not made explicit in the cited surveys, and it is one of the
more transferable findings of the study.

Finally, the collapse of the symmetric time-based self-play variant to a fail-fast
equilibrium (Appendix~\ref{app:failfast}) is best read through the game-dynamics
analysis of \cite{Letcher2019GameMechanics}. In the symmetric reward, the
goal-solving agent is penalised for the very solve time that rewards the
goal-proposing agent. This defines a near-zero-sum game whose only stable point
under the observed physics is immediate failure. The asymmetric reward of
\cite{Sukhbaatar2018IntrinsicMotivation} incentivises the proposer without
penalising the solver, and so avoids this degenerate equilibrium. The present
study provides a concrete manipulation-domain instance of why that asymmetry is
not incidental but necessary.

\subsection{Contributions Revisited}
\label{sec:contributions}

The contributions stated in Chapter~\ref{sec:introduction} can now be given
quantitative content, one per research question.

For RQ1, the thesis formulates a potential-based reward for multi-objective planar
pushing. It combines bounded exponential position and orientation potentials with
a unified $\mathrm{SE}(2)$ distance. This reward reaches 80.0\,\% scene success at
528 environments, matching the improvement reward at 2{,}048 environments and
improving the combined-objective rate by 35 percentage points.

For RQ2, the thesis provides a controlled comparison of four training regimes
under an identical reward, budget, and evaluation protocol. The forced curriculum
is statistically equivalent to the single-agent model, at 76.7\,\% against
80.0\,\%, a difference that is not significant. Both self-play variants are far
worse, at 16.7\,\% and 6.7\,\%. The cross-environment comparison establishes that
the value of the curriculum depends on the per-update batch size.

For RQ3, the thesis diagnoses the self-play failure by separating per-axis
competence from combined success. Self-play learns each of position and
orientation to roughly 50\,\%, but combines them at about 10\,\%, roughly 2.5 times
below the independence product. The goal-proposing agent, meanwhile, remains a
reliable source of valid goals. This locates the failure in skill composition
under a non-stationary goal distribution, not in skill acquisition.

\subsection{Limitations}
\label{sec:limitations}

Several limitations bound the scope of these conclusions. Validation uses a single
seed per model, so the reported scene success rates carry no confidence intervals.
Only the training-time success rate has variance data, across five chains for the
single-agent model. The best-of-20 protocol raises the scene success rate above
the trial success rate by 14 percentage points for the strongest model. Both
numbers are reported, but the headline figure is the more favourable of the two by
construction. The comparison between self-play and single-agent training differs on
more than one axis at once. These include the number of agents, the goal
distribution, the episode budget, the goal encoder, and the behavioural-cloning
term. The comparison therefore isolates self-play as implemented, rather than the
goal-proposal mechanism alone, and the axes of variation are enumerated in
Appendix~J. The sensitivity of the potential-based reward to its parameters $k_p$,
$k_r$, and $w$ was not systematically tested; the values used were chosen from
pilot runs. Finally, the off-policy SAC with Hindsight Experience Replay baseline
did not converge within the available budget, so the external calibration relies
on published planar-pushing results rather than a converged in-house baseline.

\subsection{Future Work}
\label{sec:future_work}

Each limitation suggests a concrete continuation. Repeating the validation across
multiple seeds would attach confidence intervals to the scene success rates. It
would also place the 3.3-percentage-point single-agent-against-curriculum
difference on a firmer statistical footing. A relaxed or partial-credit success
gate, which rewards satisfying each axis rather than only the conjunction, would
test whether the composition failure can be alleviated by softening the strict
gate. The symmetric time-based self-play variant collapses to a fail-fast
equilibrium (Appendix~\ref{app:failfast}); the zero-sum stability analysis of
\cite{Letcher2019GameMechanics} suggests where a stabilised formulation might be
sought. A systematic sweep of the per-update batch size would map the boundary of
the regime in which the curriculum substitutes for the shaped reward, identified
qualitatively in Section~\ref{sec:cross_env_results}. Completing the SAC with
Hindsight Experience Replay baseline to convergence would supply the off-policy
calibration that the present study leaves open. Finally, transferring the shaped
single-agent policy to a physical UR5e would test whether the efficiency gain
survives the reality gap.

Beyond these direct continuations, the relation to prior work in
Section~\ref{sec:relation} points to two deeper questions. The first is whether
self-play can be made competitive by addressing the specific failure identified,
rather than by adding scale. The value term of the Generalized Advantage
Estimation advantage is biased by the non-stationary goal distribution
(Section~\ref{sec:pbrs_asp}). A goal-proposing agent that changes the distribution
more slowly might keep the value estimate accurate enough for the two skills to
compose. One could constrain the rate of change of the proposed goal distribution,
or train the goal-solving agent against a mixture of current and past proposers.
This would test directly whether the composition failure is intrinsic to
self-play or an artefact of an over-aggressive proposer. The second question
concerns the behavioural-cloning term. The clipping mechanism of
\cite{OpenAI2021AsymmetricSF} removes the imitation signal in the discrete action
space used here. An unclipped or adaptively clipped cloning loss, or one defined
on the decoded continuous push parameters rather than the discrete bins, might
restore the demonstration signal the original formulation relies on. Both
directions treat the negative self-play result not as a dead end, but as a
localisation of where the method breaks in this domain.

\subsection{Conclusion}
\label{sec:final}

This thesis compared reward shaping against structured training dynamics on a
single contact-rich planar-pushing task, under a common budget and a common
success criterion. Potential-based reward shaping reduced the parallel-environment
requirement by a factor of about four and improved combined-objective success by
35 percentage points. A forced curriculum added no value at a sufficient batch
size, and asymmetric self-play performed far worse than single-agent training. The
self-play failure was traced to an inability to combine two individually-learned
skills under a non-stationary goal distribution, not to any failure to learn the
skills themselves. The practical recommendation that follows is simple: invest
first in a principled dense reward, and add structured training dynamics only in
the small-batch regime where they demonstrably contribute.
